\subsection{A joint normalization test using two printed VIII multiplication constants}

This supplement distinguishes a change of the formal source
parameter from an additional grading-to-Frobenius comparison.
It also retains the distinction between an ambient transition
coefficient and a canonical multiplication constant.

\paragraph{The explicit source convention.}
Write \(u\) for IV's parameter in
\[
                     (T_s-u)(T_s+u^{-1})=0.
\]
IV, \S3, equation~(22), p.~20, says this Hecke parameter agrees
with the usual quantum-group parameter; \S3.1, p.~21, after
equation~(24), explicitly gives \(q_{\mathrm{KL}}^{1/2}=u\).
VIII, \S3, p.~18, immediately after equation~(11), says that
the role of \(q\) in its \S2.2 is played by \(q^{1/2}\)
in the right-hand Hecke section \cite{IV,VIII}.
Thus, writing \(q_{\mathrm{src}}\) for that latter parameter,
the matching convention is
\begin{equation}\label{eq:jointSourceParameter}
                          q_{\mathrm{src}}=u^2.
\end{equation}
Neither source statement assigns a finite-field Frobenius
eigenvalue to the formal parameter.

There is an independent check in the actual presentations.
With \(D=(u+u^{-1})^2=q_{\mathrm{src}}+2+q_{\mathrm{src}}^{-1}\),
the displayed scalars of VIII \S5.1, p.~28, satisfy
\[
 b_2=D,\qquad -b_1=(D-2)^2,\qquad
 c_1=D(D-2)^2,\qquad c_2=D+(D-2)^2.
\]
Indeed VIII prints
\[
\begin{aligned}
 b_1&=-(q^2+1)^2/q^2,& b_2&=(q+1)^2/q,\\
 c_1&=(q^6+2q^5+3q^4+4q^3+3q^2+2q+1)/q^3,&
 c_2&=(q^4+q^3+4q^2+q+1)/q^2.
\end{aligned}
\]
These identities follow by expansion and agree with IV
Theorem~11.16, p.~65, equations~(111)--(112):
the two roots are \(a_1^2=D\) and \(a_2^2=(D-2)^2\)
\cite{IV,VIII}.
The repeated printed \(b_1\) in VIII equation~(17) is not
used here to distinguish the two basis families.

\paragraph{The three exact coefficients.}
VIII Figure~1, p.~30, label \texttt{fw1w2}, contains
\begin{equation}\label{eq:jointTransitionCoefficient}
 h(q)=4(1+q)/\sqrt q=4(q^{1/2}+q^{-1/2}).
\end{equation}
It occurs in rows~13 and~18, columns~1 and~3.
The caption and \S5.1.2, p.~29, identify the table as
ambient transition coefficients in a symmetrized
Kazhdan--Lusztig basis of \(H_3(q)\otimes H_3(q)\),
not as multiplication constants.
There is a printed column-label inconsistency:
the caption names \(W_1,W_2\), but the next explanatory
sentence calls the first columns \(\gamma_1^1,\gamma_{12}^1\),
which the preceding cell list places in \(V_1\).
The following arguments need only the unambiguous raw
entry and its transition role; they do not resolve that
label inconsistency.

The same example's Figure~2, p.~34, label
\texttt{fmultkron2}, gives two actual multiplication constants:
\begin{align}
 c(q)&=c(\gamma_2^1,\gamma_1^1;\gamma_{21}^1)
   =-\frac{(q^2+q+1)(q-1)^2}{q^2}
   =-q^2+q+q^{-1}-q^{-2},
   \label{eq:jointNegativeCoefficient}\\
 d(q)&=c(\gamma_1^1,\gamma_1^1;\Sigma)
   =\frac{(1+q)^2}{q}
   =q+2+q^{-1}.
   \label{eq:jointPositiveCoefficient}
\end{align}
The latter is printed equal to three further constants
with output \(\Sigma\); the displayed triple suffices.
Section~5.1.2 explicitly distinguishes the transition
table from these signed, vanishing-factor multiplication
constants \cite{VIII}.
In the parameter of IV,
\[
                c(u^2)=-(u-u^{-1})^2(u^2+1+u^{-2}).
\]

Fix a prime power \(Q>1\), a prime \(\ell\nmid Q\),
and a finite \(\ell\)-adic coefficient field large enough
for the eigenvalues used. We use geometric Frobenius
and the positive complex values of rational powers of \(Q\).
The integral Deligne weight convention is explicit:
an eigenvalue of weight \(m\in\mathbb Z\) has absolute
value \(Q^{m/2}\).

\begin{finding}[\fid{F65}, joint raw-coefficient normalization]
\label{f:jointRawCoefficientNormalization}
Let \(a\in\mathbb Q\setminus\{0\}\), and use the common
all-extension parameter rule \(q_{\mathrm{src}}=Q^{an}\).

First, the following two hypotheses are incompatible:
\(h(Q^{an})\) is the Euler trace of a bounded mixed point
complex with integral weights, and \(c(Q^{an})\) is
the Euler trace of a bounded pure point complex of one
integral weight.

Second, even without any ambient-transition comparison,
the two actual multiplication constants
\(d(Q^{an})\) and \(c(Q^{an})\) cannot both be Euler
traces of bounded pure point complexes of integral weights.
Their pure complex weights are allowed to be different.
In both statements, equality for all sufficiently large
extension degrees is enough.
\end{finding}

\begin{proof}
For distinct nonzero complex numbers \(\beta_1,\ldots,\beta_s\),
the sequences \(\beta_j^n\) are linearly independent,
even for all sufficiently large \(n\): take \(s\)
consecutive such degrees and use the Vandermonde matrix
times an invertible diagonal matrix.
Thus all-extension Euler traces determine the grouped
multiplicity of every visible root.

A mixed integral-weight complex may have cancellations,
but it cannot supply a visible root absent from its
cohomology. Every visible root must consequently have
integer weight.
For a pure point complex of weight \(w\), BBD~5.1.8,
pp.~126--127, puts a root of point weight \(m\)
only in degree \(m-w\) \cite{BBD}.
It cannot cancel across different degrees, and its
nonzero Euler coefficient has sign \((-1)^{m-w}\).
This argument allows Jordan blocks.

For the first statement,
\[
             h(Q^{an})=4(Q^{a/2})^n+4(Q^{-a/2})^n.
\]
The two roots are distinct since \(a\ne0\). Their weights
are \(a,-a\), so the mixed integral-weight hypothesis
forces \(a\in\mathbb Z\). But
\[
 c(Q^{an})=-(Q^{2a})^n+(Q^a)^n
                  +(Q^{-a})^n-(Q^{-2a})^n
\]
has four distinct roots of even integer weights
\(4a,2a,-2a,-4a\), and coefficients \(-1,1,1,-1\).
In any pure complex of weight \(w\), the permitted
degrees all have parity \(w\); all visible coefficients
would have the same sign. This is impossible.

For the second statement, the positive coefficient has
trace
\[
                     d(Q^{an})=(Q^a)^n+2+(Q^{-a})^n.
\]
If its pure complex weight is \(w_d\), the root \(1\)
can occur only in degree \(-w_d\). Its positive Euler
multiplicity \(2\) forces \(w_d\) even.
The other two roots have weights \(2a,-2a\); their
positive Euler multiplicities force
\(2a-w_d\) and \(-2a-w_d\) to be even integers.
Therefore \(a\) is an integer. The preceding
mixed-sign argument for \(c\) now gives the contradiction,
independently of the pure weight assigned to its complex.

Additional invisible eigenvalues cannot repair either
pure-complex contradiction: each eigenvalue in such a
complex has only one permitted cohomological degree.
\end{proof}

The first statement explicitly uses two different coefficient
types. The second needs only the two printed multiplication
constants, but still adds a pure point Euler comparison
that VIII's signed/factored positivity does not assert.
The case \(a=0\), giving \(h(1)=8\) and \(c(1)=0\),
is deliberately excluded.

A claimed realization of both multiplication constants
by pushforwards of pure objects along proper structure
morphisms \(X\to\operatorname{Spec}\mathbb F_Q\),
with \(X\) of finite type, would also be excluded under
the same parameter rule. BBD~5.1.13--5.1.14, p.~128,
gives the two weight bounds; \(f_!=f_*\) and dual
stability make the point pushforward pure.
For a proper map to a positive-dimensional base,
arbitrary stalks of its pure pushforward need not be
pure. That situation needs a separate pure-point
extraction hypothesis and is not included.

\paragraph{The isolated half-step escape remains valid.}
For \(c\) alone, take \(a=1/2\) and \(v=Q^{n/2}\).
Then
\[
 c(v)=-v^2+v+v^{-1}-v^{-2},\qquad
 -c(-v)=v^2+v+v^{-1}+v^{-2}.
\]
Let \(L_e\) be the arithmetic Weil line with geometric
Frobenius eigenvalue \(Q^{e/2}\), using the positive
square root under the chosen complex embedding.
The complex
\begin{equation}\label{eq:jointIsolatedPureEscape}
 L_2[-1]\oplus L_1\oplus L_{-1}[2]\oplus L_{-2}[3]
\end{equation}
has pure weight \(1\).
Its cohomological degrees \(1,0,-2,-3\) have point
weights \(2,1,-1,-2\), exactly \(1+i\), and its
Euler trace is \(c(Q^{n/2})\).
The eigenvalues are \(\ell\)-adic units and may also
be used in continuous arithmetic Galois representations.
No geometric half-Tate motive is asserted.

The same normalization sends the transition coefficient
to \(4Q^{n/4}+4Q^{-n/4}\), whose roots have nonintegral
weights \(1/2,-1/2\), excluding an ordinary mixed
integral-weight Euler realization.
It also sends \(d\) to \(Q^{n/2}+2+Q^{-n/2}\);
the positive root \(1\) requires an even pure complex
weight, while the positive roots of weights \(1,-1\)
require an odd one. Thus the isolated pure realization
does not extend to even the second multiplication constant.

Moreover \eqref{eq:jointSourceParameter} makes this choice
\(u=Q^{n/4}\), not \(u=Q^{n/2}\).
Finite Weil representations of nonintegral weights can
be chosen; they lie outside the integral Deligne weight
hypothesis. Neither statement is a variable-independent
exclusion of all coefficient objects or all mixed
nonproper realizations.

% Separate paragraph for Part III.
\paragraph{The normalization test in Grothendieck's language.}
The formal source parameter is a split-torus character.
SGA~3, I, 4.7.2--3 and cote~90, scan~17 provide its
action vocabulary; SGA~6, V, \S\S2--3 and cote~133,
scan~18 provide coefficient representation and
effective/virtual language \cite{SGA3,SGA6,Archive}.
They do not assign that character a Frobenius weight.
The first test crosses two different source coefficient
tables; the second avoids that issue but still requires
pure Euler interpretations of both raw multiplication constants.

The finite point complexes fit SGA~6, I, Definition~4.7,
pp.~121--122 and Example~4.8, p.~122.
The integral weights and degree-shifted purity are
BBD~5.1.5 and~5.1.8, pp.~125--127.
The proper structure-morphism consequence uses
BBD~5.1.13--5.1.14, p.~128, with \(f_!=f_*\) and dual
stability; no purity of arbitrary stalks is inferred
\cite{SGA6,BBD}.
Vandermonde uniqueness and the parity deductions are
proved above. No archival text supplies the additional
arithmetic comparisons, a canonical projector, or a
\(\lambda\)-ring structure on the noncommutative word algebra.

